% ===================== 第 9 章 =====================
\section{可靠性工程：状态机、续跑与降级}

\subsection{显式状态机：跳转非法即报错}

\code{packages/epistemo/state\_machine.py} 的 \code{TaskStateMachine} 冻结全部合法状态与转移，阶段推进只允许走到下一阶段，跳步抛 \code{InvalidTransition}。任务主状态机如图~\ref{fig:states}。

\begin{figure}[ht]
\centering
\resizebox{\linewidth}{!}{%
\begin{tikzpicture}[node distance=2.2mm and 5mm]
  \node[nd] (a) {待确认主张};
  \node[nd, right=of a] (q) {QUEUED};
  \node[nda, right=of q] (r) {RUNNING};
  \node[nd, below=7mm of q] (w) {等待人类引导\newline\tiny AWAITING\_COUNCIL\_INPUT};
  \node[ndd, right=12mm of r] (t1) {COMPLETED};
  \node[nd, below=3mm of t1] (t2) {COMPLETED\_WITH\_GAPS};
  \node[nd, below=3mm of t2] (t3) {FAILED / CANCELLED};
  \draw[flow] (a)--(q); \draw[flow] (q)--(r);
  \draw[flow] (r)--(w); \draw[flow] (w)--(q);
  \draw[flow] (r)--(t1); \draw[flow] (r)--(t2); \draw[flow] (r)--(t3);
\end{tikzpicture}}
\caption{任务主状态机：人类引导检查点回到队列续跑；终态三分支区分完整完成、带缺口完成与失败/取消。}
\label{fig:states}
\end{figure}

\subsection{事务语义：一次议会跑在一个未提交事务里}

Worker 的 \code{deliberate()} 把八阶段议会放在一次数据库事务内执行，提交时账本与证据图同时生效，不存在半截持久状态。这一选择直接决定了「暂停」的语义：

\begin{itemize}[leftmargin=1.6em, itemsep=3pt]
\item \textbf{队列级暂停}：任务在 \code{QUEUED} 与 \code{PAUSED} 之间切换，暂停期间任何 Worker 都认领不到，恢复后继续，端到端测试覆盖此保证。
\item \textbf{议会内检查点}：联合建模前唯一的固定暂停点，承载人类方向性引导（第 4.5 节）。
\end{itemize}

\subsection{检查点与两种「重新研究」}

\code{CouncilCheckpoint}（JSONB）保存已跑阶段、携带数据、未填槽位、失败记录、各阶段快照与续跑起点；\code{first\_unfinished\_phase()} 把「从未跑到」的阶段也算作未完成，因此任务取消在哪一步，就从下一步原位续跑。终态任务同样保留检查点，「从断点处研究」能还原协议实际执行轨迹。重新研究分两种模式，区别是原则性的：

\begin{mechbox}
\textbf{从头研究（full）}：克隆为全新任务（新 id、新账本），仅继承已确认主张，从预承诺重新开始；原任务保留为审计历史。必须克隆而非原位重跑——同一任务重跑会与既有幂等键冲突，上传的 PDF 等对象通过 \code{rehome\_uploaded\_objects} 迁移到新任务。\\[4pt]
\textbf{从断点研究（first\_gap）}：从第一个未完成阶段原位续跑，账本幂等性保证已完成阶段重放为空操作。
\end{mechbox}

\subsection{单席失败降级与停止即时生效}

单席模型调用失败经 \code{UnavailableDeliverator} 转为 \code{SEAT\_UNAVAILABLE} 事件（带原因与重试次数），重试至多 2 次、整轮 15 分钟封顶后记缺席，其余席位继续，缺席者在报告局限中逐条点名。研究者点击停止时，请求写入侧信道表，Worker 在阶段间以 1 秒粒度轮询并抛 \code{CouncilCancelled}，任务转入 \code{CANCELLED} 且保留全部检查点作为续跑素材。

\subsection{实时流：两条 SSE，一份一致性契约}

工作台靠两条 Server-Sent Events 流驱动：账本流（\code{/api/stream/\{task\}}）承载正式事件并支持 \code{Last-Event-ID} 续传；过程流（\code{/api/stream/\{task\}/process}）承载席位思考、token 与检索命中这类易逝数据，重连时重放尾部并按序号去重。SSE 帧只使用 \code{id:} 与 \code{data:}、事件类型放在 body 的 \code{kind} 字段——类型化帧会让未穷举词汇表的客户端静默丢事件，审计轨迹曾因此只显示 53 个事件中的 37 个。

打开被反复续跑过的历史任务时，两条流都可能重放出数千帧。客户端因此做了三层有界化：账本帧按 250 毫秒批量合并，一次突发只做一次排序与渲染；账本与过程流分别设 4000/5000 条内存上限并保留结构锚点；首次载入设 30 秒超时看门狗，超时即落到可操作的错误面板（重试 / 返回列表），杜绝「载入中」永久死机。过程流在服务端缓冲批写，且其表不设外键——Worker 持有任务行锁直到事务提交，外键的 KEY SHARE 锁会与之死锁，这是用一次整个 Worker 冻结的事故换来的设计。

\subsection{会话删除与三层测试纪律}

\code{DELETE /api/tasks/\{id\}} 物理删除任务及全部子记录，这是权限模型唯一的例外（迁移 0019 单独授予应用角色 DELETE），界面提供两段式确认与「清空全部」。测试分三层：无外部依赖的单元测试、必须真实连接 PostgreSQL 的集成测试（无数据库时显式 skip 并说明，绝不静默全绿）、Golden 因果蕴含用例与发布前 E2E 门禁。当前共 \textbf{1187 项自动化测试}；业务规则、状态机与图约束遵循先测后写，修复缺陷必须附带回归测试。
